\chapter{Known limitations}\label{chapter-16}
These findings describe the documented application implementation. They are implementation constraints, not speculative vulnerabilities or a list of completed fixes.

\Needspace{10\baselineskip}
\section*{Runtime and topology}
\pdfbookmark[1]{Runtime and topology}{chapter-16-s1}
\begingroup
\fontsize{8.7}{11.8}\selectfont
\setlength{\tabcolsep}{5pt}
\renewcommand{\arraystretch}{1.23}
\begin{longtable}{@{}>{\raggedright\arraybackslash}p{\dimexpr 0.10000\linewidth-2.000pt\relax}>{\raggedright\arraybackslash}p{\dimexpr 0.45000\linewidth-9.000pt\relax}>{\raggedright\arraybackslash}p{\dimexpr 0.45000\linewidth-9.000pt\relax}@{}}
\rowcolor{navy} \color{white}\bfseries ID & \color{white}\bfseries Finding & \color{white}\bfseries Practical consequence \\
\endfirsthead
\rowcolor{navy} \color{white}\bfseries ID & \color{white}\bfseries Finding & \color{white}\bfseries Practical consequence \\
\endhead
\endfoot
\rowcolor{pale}
GS-01 & Desktop creates one PMU/proxy from the first PMU and Local PDC. & Multiple nodes/links do not create simultaneous independent flows. \\
GS-02 & Runtime fixes loopback PMU input, UDP, 30 frames/s, 50 Hz, 120 V, and PMU ID 1. & PMU IP, reporting rate, ID and node transport edits do not control those settings. \\
\rowcolor{pale}
GS-03 & Threat Agent attachments/scope are not consulted by the runtime engine. & Attacks can run without an attached agent; selected-link scope is not enforcement. \\
GS-04 & Saved agent attack metadata is not loaded into the engine. & Loading a topology does not arm its saved attack policy. \\
\rowcolor{pale}
GS-05 & Engine counters/state/schedule survive Stop/Run. & Reset explicitly for comparable runs; a finite window may already be over. \\
GS-06 & Ordinary Open can replace a running diagram; scenario selection can still alter attacks when demo loading is blocked. & Stop before loading diagrams or scenarios. \\
\rowcolor{pale}
GS-07 & Top-level Attacks menu ignores its selected module argument. & Use the Attacks-tab dropdown and enable action. \\
\end{longtable}
\endgroup

{\footnotesize\color{muted}\raggedright Sources: \href{https://github.com/neerajjayesh/GridSec-project/blob/main/gui/main_window.py}{MainWindow} (\code{\seqsplit{gui/main\_window.py}}), \href{https://github.com/neerajjayesh/GridSec-project/blob/main/gui/attack_panel.py}{AttackPanel} (\code{\seqsplit{gui/attack\_panel.py}}), \href{https://github.com/neerajjayesh/GridSec-project/blob/main/gui/canvas.py}{Canvas} (\code{\seqsplit{gui/canvas.py}}).\par}

\Needspace{10\baselineskip}
\section*{Packet integrity and observation}
\pdfbookmark[1]{Packet integrity and observation}{chapter-16-s2}
\begingroup
\fontsize{8.7}{11.8}\selectfont
\setlength{\tabcolsep}{5pt}
\renewcommand{\arraystretch}{1.23}
\begin{longtable}{@{}>{\raggedright\arraybackslash}p{\dimexpr 0.10000\linewidth-2.000pt\relax}>{\raggedright\arraybackslash}p{\dimexpr 0.45000\linewidth-9.000pt\relax}>{\raggedright\arraybackslash}p{\dimexpr 0.45000\linewidth-9.000pt\relax}@{}}
\rowcolor{navy} \color{white}\bfseries ID & \color{white}\bfseries Finding & \color{white}\bfseries Practical consequence \\
\endfirsthead
\rowcolor{navy} \color{white}\bfseries ID & \color{white}\bfseries Finding & \color{white}\bfseries Practical consequence \\
\endhead
\endfoot
\rowcolor{pale}
GS-08 & Parser global codec has zero digital words; PMU generates one; proxy does not call \code{\seqsplit{set\_codec()}}. & Modified frames can omit the digital word and no longer match initial CFG-2 layout. \\
GS-09 & Rebuild exceptions fall back to original raw bytes; modified classification remains. & An attacked record alone does not prove changed outgoing bytes. \\
\rowcolor{pale}
GS-10 & Records precede forwarding and plots accept dropped/non-data records. & Plots and dashboard do not verify receiver delivery; startup artifacts and loss without gaps are possible. \\
GS-11 & Delay returns no content-modified/drop flag. & Delayed packets appear clean and do not create attacked SSE/syslog events. \\
\rowcolor{pale}
GS-12 & Auxiliary protocols use original measurements and independent singleton simulators. & C37.118 attack effects do not propagate to GOOSE/DNP3/Modbus. \\
\end{longtable}
\endgroup

{\footnotesize\color{muted}\raggedright Sources: \href{https://github.com/neerajjayesh/GridSec-project/blob/main/core/packet_parser.py}{Parser} (\code{\seqsplit{core/packet\_parser.py}}), \href{https://github.com/neerajjayesh/GridSec-project/blob/main/core/pdc_proxy.py}{Proxy} (\code{\seqsplit{core/pdc\_proxy.py}}), \href{https://github.com/neerajjayesh/GridSec-project/blob/main/core/attack_engine.py}{Engine} (\code{\seqsplit{core/attack\_engine.py}}), \href{https://github.com/neerajjayesh/GridSec-project/blob/main/gui/waveform_viewer.py}{Viewer} (\code{\seqsplit{gui/waveform\_viewer.py}}).\par}

\Needspace{10\baselineskip}
\section*{Integration and API}
\pdfbookmark[1]{Integration and API}{chapter-16-s3}
\begingroup
\fontsize{8.7}{11.8}\selectfont
\setlength{\tabcolsep}{5pt}
\renewcommand{\arraystretch}{1.23}
\begin{longtable}{@{}>{\raggedright\arraybackslash}p{\dimexpr 0.10000\linewidth-2.000pt\relax}>{\raggedright\arraybackslash}p{\dimexpr 0.45000\linewidth-9.000pt\relax}>{\raggedright\arraybackslash}p{\dimexpr 0.45000\linewidth-9.000pt\relax}@{}}
\rowcolor{navy} \color{white}\bfseries ID & \color{white}\bfseries Finding & \color{white}\bfseries Practical consequence \\
\endfirsthead
\rowcolor{navy} \color{white}\bfseries ID & \color{white}\bfseries Finding & \color{white}\bfseries Practical consequence \\
\endhead
\endfoot
\rowcolor{pale}
GS-13 & Manager lacks panel-used convenience methods: \code{\seqsplit{start\_pcap\_only}}, \code{\seqsplit{stop\_pcap\_only}}, \code{\seqsplit{start\_syslog\_only}}, \code{\seqsplit{stop\_syslog\_only}}, \code{\seqsplit{test\_syslog}}, \code{\seqsplit{start\_rest\_only}}, \code{\seqsplit{stop\_rest\_only}}, \code{\seqsplit{disable\_rest\_auth}}. & Affected integration buttons can raise AttributeError. Core configure/start/stop interfaces remain available. \\
GS-14 & GUI reads missing \code{\seqsplit{PacketRecord.raw\_bytes}}. & PCAP receives empty payloads; recent packet lengths are zero. Fixing only Start Capture does not fix capture data. \\
\rowcolor{pale}
GS-15 & Desktop records synthetic addresses/ports and only C37.118 callbacks. & Exports do not faithfully represent live endpoint addresses or all auxiliary traffic. \\
GS-16 & Desktop does not set remote attack hooks. & Enable can return false; disable can report success without changing the engine. \\
\rowcolor{pale}
GS-17 & API \code{\seqsplit{proto}} filtering reads \code{\seqsplit{proto}} while incidents contain \code{\seqsplit{protocol}}. & Standard-manager filtered incident queries return empty results. \\
GS-18 & REST uses serial HTTPServer with long-lived SSE. & SSE can block other requests and delay server shutdown. \\
\rowcolor{pale}
GS-19 & Root URL is stripped to an empty path. & Use \code{\seqsplit{/api/v1/docs}}; \code{\seqsplit{/}} is not a working help route. \\
GS-20 & API state/auth are process-global; limits and POST payloads lack consistent validation. & Independent server instances are not isolated; clients can receive 500s or broken connections for bad input. \\
\rowcolor{pale}
GS-21 & Manager \code{\seqsplit{sim\_running}} is not reset by desktop Stop; attack snapshot is counters only. & API status can be stale and \code{\seqsplit{/attacks/current}} is not a complete active-policy description. \\
GS-22 & No explicit integration-manager stop in MainWindow close; PCAP open failures may be stored without raising. & Explicit lifecycle and readiness checks are needed in embedded integrations. \\
\rowcolor{pale}
GS-23 & STIX uses custom fields, static mappings, and limited samples. & Validate/adapt before importing; do not treat it as a conformance-certified or complete evidence package. \\
\end{longtable}
\endgroup

{\footnotesize\color{muted}\raggedright Sources: \href{https://github.com/neerajjayesh/GridSec-project/blob/main/gui/integration_panel.py}{IntegrationPanel} (\code{\seqsplit{gui/integration\_panel.py}}), \href{https://github.com/neerajjayesh/GridSec-project/blob/main/core/integration_manager.py}{IntegrationManager} (\code{\seqsplit{core/integration\_manager.py}}), \href{https://github.com/neerajjayesh/GridSec-project/blob/main/core/rest_api.py}{REST server} (\code{\seqsplit{core/rest\_api.py}}), \href{https://github.com/neerajjayesh/GridSec-project/blob/main/core/pcap_writer.py}{PCAP writer} (\code{\seqsplit{core/pcap\_writer.py}}).\par}

\Needspace{10\baselineskip}
\section*{Platform and project maturity}
\pdfbookmark[1]{Platform and project maturity}{chapter-16-s4}
\begin{itemize}
\item GOOSE uses Linux \code{\seqsplit{AF\_PACKET}}; native Windows live publishing is unavailable and its worker may raise an unsupported-family exception.
\item GUI startup's "running" messages do not replace socket readiness checks; startup uses fixed sleeps and several workers initialize asynchronously.
\item No multi-platform live-network interoperability suite, protocol certification, production benchmark, load test, or packaged installer is established here.
\item Existing \code{\seqsplit{test\_launch.sh}} and deployment scripts contain workstation-specific paths and copy/overwrite source. They are not portable test/release commands.
\item No dependency lockfile, CI workflow, durable incident store, TLS, account/role system, or authoritative \code{\seqsplit{LICENSE}} file is included.
\end{itemize}

See \hyperref[chapter-15]{Testing} for what has actually been verified. A future fix should update this page, affected workflow documentation, and the corresponding validation evidence together.

